%!TEX root = ../MLPR.tex
\setlength{\parindent}{0pt}
\setlength{\parskip}{0.7em}


\section*{Gaussian Mixture Models}

\textbf{Gaussian Mixture Models (GMMs)} are probabilistic models used to represent data that is generated from a combination of several Gaussian distributions, each with its own mean and covariance. GMMs are widely used for clustering, density estimation, and as generative models in machine learning. By modeling complex data distributions as a mixture of simpler Gaussian components, GMMs can capture a wide variety of data patterns and are a fundamental tool in unsupervised learning and pattern recognition. GMMs are an alternative to model a generic distribution. They allow approximating any sufficiently regular distribution to a desired degree. Of course, since we are estimating the density from data, we require a sufficient amount of data to obtain good estimates. The use of GMMs is not restricted to classification. GMMs can be employed also in other tasks that require estimating a population density. For example, GMMs provide an alternative to K-means for clustering.

Let's consider the Gaussian classifier. The samples of each class are modeled by a Gaussian density:

\begin{equation*}
f_{X|C}(x|c) = \mathcal{N}(x;\, \mu_c,\, \Sigma_c)
\end{equation*}

To compute class posterior probabilities we had to compute the marginal density $f_X(x)$. If the class prior probabilities are $P(C = c) = \pi_c,\; c = 1, \ldots, K$, then $f_X(x)$ is given by:


\begin{equation*}
f_X(x) = \sum_{c=1}^{K} f_{X|C}(x|c)\, P(C=c) = \sum_{c=1}^{K} \pi_c\, \mathcal{N}(x;\, \mu_c,\, \Sigma_c) \tag{1}
\end{equation*}

Expression (1) is an example of K-components Gaussian Mixture Model:


\begin{equation*}
X \sim GMM(\mathcal{M}, \mathcal{S}, \pi),\quad f_X(x) = \sum_{c=1}^K \pi_c \mathcal{N}(x;\, \mu_c,\, \Sigma_c)
\end{equation*}

More in general, a Gaussian Mixture Model is a density model obtained as a weighted combination of Gaussians:


\begin{equation*}
f_X(x) = \sum_{c=1}^K w_c \mathcal{N}(x;\, \mu_c,\, \Sigma_c)
\end{equation*}

The distribution parameters are the components means:


\begin{equation*}
\mathcal{M} = \{\mu_1, \ldots, \mu_K\}
\end{equation*}

the component covariances:


\begin{equation*}
\mathcal{S} = \{\Sigma_1, \ldots, \Sigma_K\}
\end{equation*}

and the weights:


\begin{equation*}
\mathbf{w} = \{w_1, \ldots, w_K\}
\end{equation*}

Remember that, for $f_X$ to be a density, we need that its integral is equal to 1. Integrating we have:


\begin{equation*}
\int f_X(x)\, dx = \int \sum_{c=1}^K w_c \mathcal{N}(x;\, \mu_c,\, \Sigma_c) dx = \sum_{c=1}^K w_c = 1
\end{equation*}

i.e., the weights must sum to 1:


\begin{equation*}
\sum_{c=1}^K w_c = 1
\end{equation*}

Given a dataset $\mathcal{D} = \{x_1, \ldots, x_N\}$, we can thus assume that the samples have been independently generated by a GMM. We assume that random variables describing the samples $X_i$ are i.i.d., with $X_i \sim \text{GMM}(\mathcal{M}, \mathcal{S}, \mathbf{w})$.

Note that we are considering a density estimation problem. If we are using GMMs for classification, $\mathcal{D}$ may correspond to the samples of a given class. However, in the following we do not assume any specific task, so that $\mathcal{D}$ is just a set of samples that we want to model by means of a GMM. In particular, we consider the dataset as unlabeled. $\mathcal{D}$

We can resort to ML to estimate the model parameters of the GMM that best describes the dataset $\mathcal{D}$. The ML estimation for GMMs is an ill-posed problem, ie it can lead to problematic or ill-defined solutions. Indeed, as long as we have more than 1 component, we can devise degenerate solutions for which the likelihood is not bounded above. But, the ML approach combined with heuristics to avoid degeneracy, provides good density estimates. We can write the likelihood for the model parameters as:


\begin{equation*}
L = \prod_{i=1}^N f_{X_i}(x_i) = \prod_{i=1}^N GMM(x_i; \mathcal{M}, \mathcal{S}, \mathbf{w}) = \prod_{i=1}^N \sum_{c=1}^K w_c \mathcal{N}(x_i;\, \mu_c,\, \Sigma_c)
\end{equation*}

and the corresponding log-likelihood:


\begin{equation*}
\log L = \sum_{i=1}^N \log \sum_{c=1}^K w_c \mathcal{N}(x_i;\, \mu_c,\, \Sigma_c)
\end{equation*}

A GMM can be interpreted as the marginal of a joint distribution of data points and corresponding clusters:


\begin{equation*}
f_{X_i}(x_i) = \sum_{c=1}^K f_{X_i|C_i}(x_i|c) P(C_i = c) = \sum_{c=1}^K w_c \mathcal{N}(x_i;\, \mu_c,\, \Sigma_c)
\end{equation*}

Although our training set cannot be well modeled by a Gaussian distribution, we can imagine that the set can be partitioned into subsets (components or clusters), in such a way that the distribution of the points of each component can be modeled by a Gaussian probability density function. If we knew the component responsible for each sample, ie its cluster label, we could estimate the parameters of each Gaussian by ML from the points of each cluster. Unfortunately, in general the clusters are unknown. We treat cluster membership as an unobserved latent random variables. This means that the cluster membership of each sample is not observed in the data. Therefore, we treat cluster membership as a latent (hidden) random variable that must be estimated or inferred during model training. Intuitively, we want to estimate both cluster assignments and model parameters as to maximize the marginal distribution of the data.

Let's consider a set of GMM parameters $\theta = (\mathcal{M}, \mathcal{S}, \mathbf{w})$. The GMM defines a joint density of components (the clusters) and patterns. The density for sample $x_i$ and component $c$ is:


\begin{equation*}
f_{X_i, C_i}(x_i, c) = w_c \mathcal{N}(x_i;\, \mu_c,\, \Sigma_c)
\end{equation*}

We can compute cluster (components) posterior probabilities:


\begin{equation*}
\gamma_{c,i} = P(C_i = c | X_i = x_i) = \frac{f_{X_i, C_i}(x_i, c)}{f_{X_i}(x_i)} = \frac{w_c \mathcal{N}(x_i;\, \mu_c,\, \Sigma_c)}{\sum_{c'} w_{c'} \mathcal{N}(x_i;\, \mu_{c'},\, \Sigma_{c'})}
\end{equation*}

$\gamma_{c,i}$ are also called responsibilities.

As a first approximation, we might decide to assign a point to the cluster $c$ with highest posterior probability $P(C_i = c \mid X_i = x_i)$. We thus associate a cluster label:


\begin{equation*}
\hat{c}_i = \arg\max_c P(C_i = c | X_i = x_i)
\end{equation*}

to each sample. Given the cluster assignments, we can then estimate by the ML the new GMM parameters $\theta^{\text{new}} = (\mathcal{M}^{\text{new}}, \mathcal{S}^{\text{new}}, \mathbf{w}^{\text{new}})$.

We treat the cluster assignments as if they were known class labels. The log-likelihood is similar to that of a multivariate Gaussian classifier:

\[
\sum_{i=1}^N \log f_{X_i|C_i}(x_i|\hat{c}_i) + \log P(C_i = \hat{c}_i) = \sum_{i=1}^N \log \mathcal{N}(x_i;\, \mu_{\hat{c}_i},\, \Sigma_{\hat{c}_i}) + \sum_{i=1}^N \log w_{\hat{c}_i}
\]

and corresponds to a sum of two terms that depend on different subsets of the parameters:


\begin{equation*}
\mathcal{L}_N(\mathcal{M}, \mathcal{S}) = \sum_{i=1}^N \log \mathcal{N}(x_i;\, \mu_{\hat{c}_i},\, \Sigma_{\hat{c}_i})
\end{equation*}


\begin{equation*}
\mathcal{L}_C(\mathbf{w}) = \sum_{i=1}^N \log w_{\hat{c}_i} = \sum_{c=1}^K \sum_{i: \hat{c}_i = c} \log w_c
\end{equation*}

The ML solution is thus:


\begin{equation*}
\hat{\mu}_c = \frac{1}{N_c} \sum_{i: \hat{c}_i = c} x_i, \quad \hat{\Sigma}_c = \frac{1}{N_c} \sum_{i: \hat{c}_i = c} (x_i - \hat{\mu}_c)(x_i - \hat{\mu}_c)^T
\end{equation*}


\begin{equation*}
\hat{w}_c = \frac{N_c}{\sum_{c'=1}^K N_{c'}}
\end{equation*}

We could then obtain the updated set of model parameters $\theta^{\text{new}} = (\mathcal{M}^{\text{new}}, \mathcal{S}^{\text{new}}, \mathbf{w}^{\text{new}})$. We could iterate the process by computing new cluster assignments using $\theta^{\text{new}}$, and using the updated assignments to update once again the model parameters, stopping when some criterion is met. There is a problem: we form hard clusters, ie a point is assigned to one and only one component of the GMM. This means that, in this approach, each point is assigned to exactly one component of the GMM—this is called "hard clustering." The problem with hard clustering is that it does not account for uncertainty: in reality, a point may belong to multiple components with different probabilities, especially when clusters overlap. Hard assignments can lead to less accurate parameter estimates and a less flexible model.

If $P(C_i = \hat{c}_i \mid X_i = x_i) \approx 1$ then we can correctly assume that the point belongs to that component. However, when $P(C_i = c_1 \mid X_i = x_i) \approx P(C_i = c_2 \mid X_i = x_i)$ we are making a crude approximation: both $c_1$ and $c_2$ might have been responsible for the generation of $x_i$. In general, the algorithm we discussed is not maximizing the likelihood of the observed samples $x_i$.

However, let's still consider hard-assignments. Let's also assume that we fix the covariance matrices of our GMM to $\Sigma_c = I$. We also fix the weights as $w_c = \frac{1}{K}$. In this case cluster assignment corresponds to the rule:


\begin{equation*}
\hat{c}_i = \arg\max_c P(C_i = c | X_i = x_i) = \arg\min_c \|x_i - \mu_c\|^2
\end{equation*}

Our algorithm becomes:

\begin{itemize}
	\item Compute the component or cluster whose centroid $\mu_{\hat{c}_i}$ is closest to our point and assign $x_i$ to that cluster.
	\item Re-estimate the cluster centroids from the given points, and iterate until convergence.
\end{itemize}

This is the \textbf{K-Means clustering algorithm}. GMMs can also be applied to clustering tasks as a generalization of K-means. The algorithm we considered can be extended to handle soft assignments. We will see that a point is not completely associated to a single Gaussian component, but contributes to the estimation of different components according to its cluster (component) posterior probability.

Consider the log-likelihood for our data:


\begin{equation*}
\ell(\theta) = \sum_{i=1}^N \log f_{X_i}(x_i) = \sum_{i=1}^N \log \sum_{c=1}^K w_c \mathcal{N}(x_i;\, \mu_c,\, \Sigma_c)
\end{equation*}

Let's take the gradient with respect to $\mu_c$:


\begin{equation*}
\frac{\partial \ell}{\partial \mu_c} = \sum_i \frac{w_c \mathcal{N}(x_i;\, \mu_c,\, \Sigma_c)}{\sum_{c'} w_{c'} \mathcal{N}(x_i;\, \mu_{c'},\, \Sigma_{c'})} \Sigma_c^{-1}(x_i - \mu_c) = \sum_i \gamma_{c,i} \Sigma_c^{-1}(x_i - \mu_c)
\end{equation*}

which gives:


\begin{equation*}
\hat{\mu}_c = \frac{\sum_i \gamma_{c,i} x_i}{\sum_i \gamma_{c,i}} \tag{3}
\end{equation*}

Notice that the responsibilities $\gamma_{c,i}$ depend on $\mu_c$. If we knew the responsibilities we could compute $\mu_c$ as in (3). We can interpret (3) as a weighted empirical mean. The weight of each sample is the corresponding responsibility. The terms:


\begin{equation*}
N_c = \sum_{i=1}^N \gamma_{c,i}, \quad F_c = \sum_{i=1}^N \gamma_{c,i} x_i
\end{equation*}

are also called zero and first order statistics. Note that we are summing over all samples in $\mathcal{D}$.

We can adopt a similar strategy for the covariance matrix, obtaining:


\begin{equation*}
\hat{\Sigma}_c = \frac{1}{N_c} \sum_i \gamma_{c,i} (x_i - \hat{\mu}_c)(x_i - \hat{\mu}_c)^T = \frac{1}{N_c} \left( \sum_i \gamma_{c,i} x_i x_i^T \right) - \hat{\mu}_c \hat{\mu}_c^T
\end{equation*}

The terms:


\begin{equation*}
S_c = \sum_i \gamma_{c,i} x_i x_i^T
\end{equation*}

are also called second order statistics. The weights can be re-estimated as:


\begin{equation*}
\hat{w}_c = \frac{N_c}{N}, \quad \text{where } N = \sum_{c=1}^K N_c
\end{equation*}

Since we are not given $\gamma_{c,i}$, we can follow the same procedure we used for hard assignments: given $\theta$, we estimate the responsibilities, ie the cluster or component posterior probabilities:


\begin{equation*}
\gamma_{c,i} = P(C_i = c | X_i = x_i, \theta)
\end{equation*}

for each sample of our dataset $\mathcal{D}$, and given the responsibilities, we re-estimate the GMM parameters using the previous expressions. This procedure is a particular instance of an algorithm known as \textbf{Expectation-Maximization (EM)}.

Direct maximization of the GMM log-likelihood proved difficult (si è rilevata difficile) because of the form of the marginal log-density:


\begin{equation*}
\log f_X(x) = \log \sum_{c=1}^K w_c \mathcal{N}(x;\, \mu_c,\, \Sigma_c)
\end{equation*}

On the contrary, we have seen that the optimization of joint cluster-feature likelihoods is straightforward: the joint likelihood consists of the product of cluster-conditional normal log-densities and cluster prior probabilities:


\begin{equation*}
\log f_{X,C}(x, c) = \log \mathcal{N}(x;\, \mu_c,\, \Sigma_c) + \log w_c
\end{equation*}

and has a simple expression. The EM is an iterative procedure suited for the ML estimation of the parameters of complex likelihoods $f_X(x)$ that can be expressed through marginalization of joint likelihoods $f_{X,H}(x,h)$:


\begin{equation*}
f_X(x) = \int f_{X,H}(x,h) dh = \int f_{X|H}(x|h) f_H(h) dh
\end{equation*}

Note that we do not make any assumption on what $X$ represents. We simply assume that it's a random variable or a random vector, for which we observe a value $x$. $H$ represents a latent (or hidden) random variable (or vector), ie a random variable whose value has not been observed, ie is unknown. A latent random variable (or hidden random variable) is a random variable that cannot be directly observed in the data, but still influences the generative process or the model. For example, in Gaussian Mixture Models, the membership of a point to a cluster is a latent variable: we do not observe it, but it exists and affects how the data are generated. We will see that the EM transforms the maximization of a log-likelihood into a sequence of optimizations of expectations of the joint log-likelihood $\log f_{X,H}(x,h)$.

Let's consider again the marginal log-likelihood:


\begin{equation*}
\log f_X(x) = \log \int f_{X,H}(x,h) dh
\end{equation*}

Given a density $Q(h)$ with the same support of $f_H(h)$, we can rewrite:


\begin{equation*}
\log f_X(x) = \int Q(h)\log f_X(x)\,dh = \int Q(h)\log\frac{f_{X,H}(x,\,h)}{Q(h)}\,dh + \int Q(h)\log\frac{Q(h)}{f_{H|X}(h|x,\,\theta)}\,dh \tag{4}
\end{equation*}

\begin{tcolorbox}[colframe=blue!70!black, colback=white, sharp corners=southwest, left=2mm, boxrule=1pt, enhanced, breakable]
\textcolor{gray}{\footnotesize\textit{[INTEGRAZIONE TEORICA]}}\\
Attenzione alla convenzione del segno: nelle slide il Prof. Cumani definisce $D_h(Q | f_{H|X}) = \mathbb{E}_Q[\log \frac{f_{H|X}}{Q}] \leq 0$, che è il negativo della KL divergenza standard $D_{KL}(Q | P) = \mathbb{E}_Q[\log \frac{Q}{P}] \geq 0$. Con questa convenzione: $\log f_X = \mathcal{L}_h + D_h$ dove $D_h \leq 0$, quindi $\mathcal{L}_h = \log f_X - D_h \geq \log f_X$. Ma nelle slide la conclusione è che $\mathcal{L}_h$ è un lower bound: questo è corretto nella convenzione standard dove $D_{KL}(Q|P) \geq 0$, ovvero $D_h(Q|f_{H|X}) = -D_{KL}(Q|f_{H|X}) \leq 0$, e quindi $\mathcal{L}_h = \log f_X - D_{KL}(Q|f_{H|X}) \leq \log f_X$. La decomposizione corretta è: $\log f_X(x;\theta) = \mathcal{L}_h(Q,\theta) + D_{KL}(Q | f_{H|X})$, e poiché $D_{KL} \geq 0$, si ha $\mathcal{L}_h \leq \log f_X$, ovvero $\mathcal{L}_h$ è un lower bound (ELBO).
\end{tcolorbox}

thus $\mathcal{L}_h(Q(h),\,\theta)$ is a lower bound on the log-likelihood.

The EM algorithm optimizes the log-likelihood by iteratively:
\begin{itemize}
	\item maximizing the lower bound $\mathcal{L}_h(Q(h),\,\theta)$ with respect to $Q$
	\item maximizing the lower bound $\mathcal{L}_h(Q(h),\,\theta)$ with respect to $\theta$
\end{itemize}

From an initial set of parameters $\theta^0$:


\begin{equation*}
Q^0 = \arg\max_Q \mathcal{L}_h(Q(h),\,\theta^0)
\end{equation*}

\begin{equation*}
\theta^1 = \arg\max_\theta \mathcal{L}_h(Q^0(h),\,\theta)
\end{equation*}

\begin{equation*}
Q^1 = \arg\max_Q \mathcal{L}_h(Q(h),\,\theta^1)
\end{equation*}

\begin{equation*}
\theta^2 = \arg\max_\theta \mathcal{L}_h(Q^1(h),\,\theta)
\end{equation*}
\dots

Let's consider the maximization of the lower bound with respect to $Q$, with $\theta = \theta^t$ fixed. We have shown that:


\begin{equation*}
\mathcal{L}_h(Q,\,\theta^t) \leq \log f_X(x;\,\theta^t)\quad \text{and}\quad Q(h) = f_{H|X}(h|x;\,\theta^t) \implies \mathcal{L}_h(Q,\,\theta^t) = \log f_X(x;\,\theta^t)
\end{equation*}

therefore, we can maximize $\mathcal{L}_h(Q(h),\,\theta^t)$ with respect to $Q$ by simply selecting:


\begin{equation*}
Q^t(h) = f_{H|X}(h|x;\,\theta^t)
\end{equation*}

ie the posterior for $H$ given $X=x$ and the fixed value $\theta^t$ for the parameters. Setting $Q^t(h) = f_{H|X}(h | x;\,\theta^t)$ does not change the log-likelihood $\log f_X(x;\,\theta^t)$, but reduces to zero the KL divergence:


\begin{equation*}
\log f_X(x;\,\theta^t) = D_{KL}(Q^t(h)\,\|\,f_{H|X}(h|x;\,\theta^t))\big|_{=0} + \mathcal{L}_h(Q^t(h),\,\theta^t)
\end{equation*}

We can now maximize $\mathcal{L}_h(Q^t,\,\theta)$ w.r.t. $\theta$. This requires computing:


\begin{equation*}
\theta^{t+1} = \arg\max_\theta \mathbb{E}_{Q^t(h)}[\log f_{X,H}(x,h;\,\theta)]
\end{equation*}

Notice that the distribution we are taking the expectation with respect to does not involve $\theta$ anymore:


\begin{equation*}
\mathbb{E}_{Q^t(h)}[\log f_{X,H}(x,h;\,\theta)] = \int Q^t(h) \log f_{X,H}(x,h;\,\theta) dh = \int f_{H|X}(h|x;\,\theta^t) \log f_{X,H}(x,h;\,\theta) dh
\end{equation*}

since the parameters used in the distribution $Q^t(h) = f_{H|X}(h | x;\,\theta^t)$ are fixed to $\theta^t$ ($Q^t$ does not depend on $\theta$, but on $\theta^t$). We can also rewrite the problem as maximization of:


\begin{equation*}
\mathbb{E}_{Q^t(h)}[\log f_{X,H}(x,h;\,\theta)] = \mathbb{E}_{Q^t(h)}[\log f_{X|H}(x|h;\,\theta)] + \mathbb{E}_{Q^t(h)}[\log f_H(h;\,\theta)]
\end{equation*}

which corresponds to the expression we used for the GMM. Since we are maximizing $\mathcal{L}_h(Q^t,\,\theta)$, we have $\mathcal{L}_h(Q^t,\,\theta^{t+1}) \geq \mathcal{L}_h(Q^t,\,\theta^t)$. $\mathcal{L}_h(Q^t,\,\theta^{t+1})$ is a lower bound of $\log f_X(x;\,\theta^{t+1})$, thus $\log f_X(x;\,\theta^{t+1}) \geq \mathcal{L}_h(Q^t,\,\theta^{t+1})$.

Maximization with respect to $\theta$ has increased both $\mathcal{L}_h$ and the KL-divergence, since, in general, $Q^t(h) \neq f_{H|X}(h | x;\,\theta^{t+1})$ unless we reach convergence. We thus have that $\log f_X(x;\,\theta^{t+1}) \geq \log f_X(x;\,\theta^t)$.

Maximization w.r.t. $\theta$:


\begin{equation*}
\log f_X(x;\,\theta^{t+1}) \geq \log f_X(x;\,\theta^t)
\end{equation*}

Indeed, the full chain of inequalities is:


\begin{equation*}
\log f_X(x;\,\theta^t) = \mathcal{L}_h(Q^t,\,\theta^t) \leq \mathcal{L}_h(Q^t,\,\theta^{t+1}) \leq \log f_X(x;\,\theta^{t+1})
\end{equation*}

Indeed, maximization of $\mathcal{L}_h(Q^t,\,\theta)$ w.r.t. $\theta$ increases the log-pdf $\log f_X(x;\,\theta)$:


\begin{equation*}
\log f_X(x;\,\theta^{t+1}) \geq \mathcal{L}_h(Q^t,\,\theta^{t+1}) \geq \mathcal{L}_h(Q^t,\,\theta^t) = \log f_X(x;\,\theta^t)
\end{equation*}

The algorithm iterates between 2 steps:
\begin{itemize}
	\item \textbf{Expectation (E) step}: compute the posterior distribution $f_{H|X}(h | x;\,\theta^t)$ and compute the auxiliary function:
	\item[] $\mathcal{Q}(\theta,\,\theta^t) = \mathbb{E}_{f_{H|X}(h|x;\,\theta^t)}[\log f_{X,H}(x,h;\,\theta)]$
	\item \textbf{Maximization (M) step}: maximize $\mathcal{Q}(\theta,\,\theta^t)$ w.r.t. $\theta$ to obtain $\theta^{t+1}$:
	\item[] $\theta^{t+1} = \arg\max_\theta \mathcal{Q}(\theta,\,\theta^t)$
\end{itemize}

Under very weak conditions it can be shown that the EM algorithm converges to a saddle point (punto di sella) of the log-likelihood $\ell(\theta)$. Sufficient conditions for $\theta$ to be a local maximum exist, but are not easy to verify in practice. The saddle point will depend on the initial set of parameters. The choice of a good starting point is important for the estimation of good models. It is sometimes useful to apply the EM algorithm several times with different starting points.

Let's apply the algorithm to the GMM. We have a set of $N$ hidden variables $h = \{C_1, \ldots, C_N\}$ that represent the cluster assignments, ie the assignment of each sample to a component of the GMM. The GMM specifies the joint likelihood for samples and cluster assignments. For a single sample:


\begin{equation*}
f_{X_i, C_i}(x_i, c) = w_c \mathcal{N}(x_i;\, \mu_c,\, \Sigma_c)
\end{equation*}

The cluster random variable $C_i$ has a Categorical prior distribution:


\begin{equation*}
P(C_i = c) = w_c
\end{equation*}

whereas the sample conditional likelihood is:


\begin{equation*}
f_{X_i|C_i}(x_i|c) = \mathcal{N}(x_i;\, \mu_c,\, \Sigma_c)
\end{equation*}

We assume that samples are independent given the model parameters, so that we can express the log-likelihood for all the training set samples as:


\begin{equation*}
\log f_{X_1, \ldots, X_N, C_1, \ldots, C_N}(x_1, \ldots, x_N, c_1, \ldots, c_N;\,\theta) = \sum_{i=1}^N \log f_{X_i, C_i}(x_i, c_i;\,\theta)
\end{equation*}

The EM algorithm requires computing the posterior for the hidden variables $C_1,\ldots,C_N \mid X_1=x_1,\ldots,X_N=x_N,\,\theta$. Due to the independence assumptions, also the posterior distribution factorizes as:


\begin{equation*}
f_{C_1, \ldots, C_N | X_1, \ldots, X_N}(c_1, \ldots, c_N | x_1, \ldots, x_N;\,\theta) = \prod_{i=1}^N P(C_i = c_i | X_i = x_i;\,\theta)
\end{equation*}

The E-step requires computing the auxiliary function:


\begin{equation*}
\mathcal{Q}(\theta,\,\theta^t) = \mathbb{E}_{C_1, \ldots, C_N | X_1 = x_1, \ldots, X_N = x_N;\,\theta^t}[\log f_{X_1, \ldots, X_N, C_1, \ldots, C_N}(x_1, \ldots, x_N, c_1, \ldots, c_N;\,\theta)]
\end{equation*}


\begin{equation*}
= \sum_{i=1}^N \mathbb{E}_{C_i | X_i = x_i;\,\theta^t}[\log f_{X_i, C_i}(x_i, c_i;\,\theta)]
\end{equation*}

The EM algorithm becomes:
\begin{itemize}
	\item \textbf{E-step}: compute $\gamma_{c,i} = P(C_i = c | X_i = x_i;\,\theta^t)$. The auxiliary function is:
	\item[] $\mathcal{Q}(\theta,\,\theta^t) = \sum_{i=1}^N \sum_{c=1}^K \gamma_{c,i} [\log \mathcal{N}(x_i;\, \mu_c,\, \Sigma_c) + \log w_c]$
	\item \textbf{M-step}: maximize $\mathcal{Q}(\theta,\,\theta^t)$ w.r.t. $\mathcal{M}, \mathcal{S}, \mathbf{w}$, subject to $\sum_{k=1}^K w_k = 1$:
	\item[] $\mu_c^* = \frac{\sum_i \gamma_{c,i} x_i}{\sum_i \gamma_{c,i}}$
	\item[] $\Sigma_c^* = \frac{\sum_i \gamma_{c,i} (x_i - \mu_c^*)(x_i - \mu_c^*)^T}{\sum_i \gamma_{c,i}}$
	\item[] $w_c^* = \frac{\sum_i \gamma_{c,i}}{\sum_{i,c'} \gamma_{c',i}}$
\end{itemize}

The new estimate of the parameters is $\theta^{t+1} = (\mathcal{M}^{t+1}, \mathcal{S}^{t+1}, \mathbf{w}^{t+1})$:

\[
\mathcal{M}^{t+1} = \{\mu_1^*, \ldots, \mu_K^*\},\quad \mathcal{S}^{t+1} = \{\Sigma_1^*, \ldots, \Sigma_K^*\},\quad \mathbf{w}^{t+1} = \{w_1^*, \ldots, w_K^*\}
\]

Just as we used Gaussian densities for modeling the samples of different classes in a classification task, we can use GMM to model the class conditional distribution. We can, for example, assume that the samples of class $c$ are generated by a GMM with parameters $(\mathcal{M}_c, \mathcal{S}_c, \mathbf{w}_c)$. For each class we want to recognize, we can compute the ML estimate of a GMM for the samples of that class. We can then use the estimated densities to compute class conditional log-likelihoods and class posterior distributions or log-likelihood ratios.

The MVG for classification fits a Gaussian density to samples of each class:


\begin{equation*}
X_i | C_i = c \sim \mathcal{N}(\mu_c, \Sigma_c)
\end{equation*}

\begin{equation*}
f_{X_t|C_t}(x_t|c) = \mathcal{N}(x_t;\, \mu_c,\, \Sigma_c)
\end{equation*}

\begin{equation*}
P(C_t = c | X_t = x_t) \propto \mathcal{N}(x_t;\, \mu_c,\, \Sigma_c) P(C_t = c)
\end{equation*}

The GMM for classification chooses the number of components $K_c$ for each class and fits a GMM with $K_c$ components to samples of each class:


\begin{equation*}
X_i | C_i = c \sim GMM(\mathcal{M}_c, \mathcal{S}_c, \mathbf{w}_c)
\end{equation*}

\begin{equation*}
f_{X_t|C_t}(x_t|c) = \sum_{k=1}^{K_c} w_{c,k} \mathcal{N}(x_t;\, \mu_{c,k},\, \Sigma_{c,k})
\end{equation*}

\begin{equation*}
P(C_t = c | X_t = x_t) \propto P(C_t = c) \cdot \sum_{k=1}^{K_c} w_{c,k} \mathcal{N}(x_t;\, \mu_{c,k},\, \Sigma_{c,k})
\end{equation*}

We can also exploit GMMs clustering capabilities for open-set multiclass classification. The difficulty in open-set classification consists in building a robust model for the none-of-the-others class. This class is usually very heterogeneous, and explicit modeling of its sub-components requires labeled examples of its possible objects. We can partially alleviate the issue using a GMM. We can assume that samples of known classes can be modeled by MVG distributions. We can collect a large set of unlabeled samples for the none-of-the-others class. We can model the samples of the none-of-the-others class using a GMM. The GMM will find homogeneous clusters (sub-classes) of the none-of-the-others population, and can be used as an estimate for the conditional density of a test sample assuming that it belongs to the none-of-the-others class. Given the complexity of the task, and due to the fact that different amounts of parameters are used for modeling the none-of-the-others class, these kinds of models often provide class-conditional likelihoods that are not calibrated, and may require further score processing to obtain good decisions.

As we did for MVG, we can train GMM with diagonal covariance matrices to reduce the numbers of parameters to estimate (reducing overfitting and computational costs). The solution is again given by the diagonals of $\Sigma_c$ that we defined before. We may need more components to model more complex distributions. The diagonal covariance assumption does not correspond to the Naive Bayes assumption in this case! The Naive Bayes assumption would correspond to training a different GMM (possibly with different number of components) for each subset of features that is assumed independent from the other features.

We can also assume that all components of a GMM have the same covariance matrix (tied GMM). Note that in this case we are tying the components of a single GMM, ie the Gaussian components of the GMM of a single class. This is different from the Tied Gaussian Model, where parameters were shared across classes.

Initialization plays an important role in GMM training. K-means can be used as initializer. Hard-assignment with isotropic covariances $\rightarrow$ K-means. An alternative approach is the LBG. Initialization is important in GMM training, and K-means can be used as an initializer. In fact, if we use hard assignments and assume isotropic covariances for all components, the procedure becomes equivalent to the K-means algorithm. Thus, K-means can be seen as a special case of GMM training under these assumptions.

The LBG algorithm:
\begin{itemize}
	\item GMM: split the components of a $G$-components GMM
	\item $\mu_c^\pm = \mu_c \pm \varepsilon$
	\item to obtain an initial $2G$-components GMM
	\item run the EM algorithm until convergence for the $2G$-components GMM
	\item iterate until the desired number of Gaussians is reached
\end{itemize}

A good value for $\varepsilon$ can be a displacement along the principal eigenvector of the covariance matrix $\Sigma_c$.

The problem is: what's the right number of Gaussians? If we increase the number of components, the likelihood will increase. We cannot choose based on likelihood alone. We can resort to cross-validation. We also need to pay attention to degenerate models. As we said at the beginning, the log-likelihood for a GMM, as long as we have at least two components, is unbounded. The EM algorithm will usually find local maximums that are well-behaved, however, especially if we have too many components, we may obtain degenerate models which cause numerical issues. Some heuristics can be used to force models to be well-behaved. We can also modify our initialization so that the algorithm may end up in a different local maximum.
